
\section{10. Limitations}

\subsection{10.1 Task Coverage}

Only \textbf{SST-2} is evaluated. Adaptive SVD may behave differently on:
\begin{itemize}
\item Generation tasks (summarization, translation)
\end{itemize}
\begin{itemize}
\item Multi-class classification (MNLI, AG News)
\end{itemize}
\begin{itemize}
\item Token-level tasks (NER, POS)
\end{itemize}
\begin{itemize}
\item Long-context tasks (>512 tokens)
\end{itemize}

\subsection{10.2 Model Coverage}

We test on three models (66.9M--124.6M params). Not tested:
\begin{itemize}
\item Large models (>1B)
\end{itemize}
\begin{itemize}
\item Deep models (>24 layers)
\end{itemize}
\begin{itemize}
\item Mixture-of-Experts
\end{itemize}
\begin{itemize}
\item Vision transformers
\end{itemize}
\begin{itemize}
\item Multimodal transformers
\end{itemize}

\subsection{10.3 Hardware}

All experiments on \textbf{RTX 4070 Laptop (8 GB)}. Behavior on:
\begin{itemize}
\item A100/H100: SVD faster, method should scale
\end{itemize}
\begin{itemize}
\item CPU inference: Two matmuls may be slower
\end{itemize}
\begin{itemize}
\item Edge devices (Jetson): Untested
\end{itemize}

\subsection{10.4 Inference Latency}

Factorized layers are slower at inference:
\begin{itemize}
\item Two matmuls instead of one
\end{itemize}
\begin{itemize}
\item Intermediate $r$-dim tensor materialized
\end{itemize}
\begin{itemize}
\item Adapter adds one more matmul
\end{itemize}
\begin{itemize}
\item ~{{LATENCY\_RATIO}}$\times$ slower on GPU
\end{itemize}

\textbf{Mitigation.} The compressed model uses less memory, so larger batch sizes 
may fit, potentially offsetting latency.

\subsection{10.5 Per-Projection Rank is Heuristic}

The energy threshold is a \textbf{principled heuristic}, not an optimization. The 
optimal rank for downstream task may differ from the energy-thresholded rank.

\subsection{10.6 No Downstream Error Bound}

We provide:
\begin{itemize}
\item \textbf{Reconstruction error bound} (Prop 4.9.1)
\end{itemize}
\begin{itemize}
\item \textbf{No bound} on downstream task accuracy
\end{itemize}

The gap between reconstruction and task accuracy is not characterized.

\subsection{10.7 Residual Adapter Overhead}

The adapter adds $r_L (d_\text{in} + d_\text{out})$ parameters per projection. 
For $d = 768$ and $r_L = 16$:

\[
\text{Overhead per projection} = 16 \cdot 1536 = 24,576 \text{ params}
\]

vs. the compressed base of $r_P \cdot 1536$. This is ~{{ADAPTER\_OVERHEAD}}\% 
of the base.

\subsection{10.8 Single-Seed Results}

All results from a single seed. Standard error $\approx 1.69\%$ at $N = 872$. 
Differences below 3.4\% may be noise.

\subsection{10.9 No Comparison with All Baselines}

We do not compare with:
\begin{itemize}
\item DistilBERT / TinyBERT (distillation)
\end{itemize}
\begin{itemize}
\item Performer / Linformer (attention replacement)
\end{itemize}
\begin{itemize}
\item SparseGPT / Wanda (pruning)
\end{itemize}
\begin{itemize}
\item Direct head pruning
\end{itemize}

\subsection{10.10 T5 and Qwen Verification}

T5 support is {{T5\_VERIFIED}}, Qwen support is {{QWEN\_VERIFIED}}. Full 
verification on these architectures is deferred to future work.

\subsection{10.11 Theoretical Guarantee on Per-Layer Optimality}

The energy threshold is not provably optimal. Alternatives (learned $\tau$, 
per-layer $\tau$, activation-aware $\tau$) are not tested.

\subsection{10.12 Code Quality}

Research-grade:
\begin{itemize}
\item Broad exception handling
\end{itemize}
\begin{itemize}
\item No unit tests for edge cases
\end{itemize}
\begin{itemize}
\item Numerical stability could be improved
\end{itemize}
\begin{itemize}
\item T5 bias path untested on decoder side
\end{itemize}

\subsection{10.13 Ethical Considerations}

Compression is dual-use. See Appendix E.